\section{Empirical Results}
\label{sec:empirical}

\subsection{Setup}

We evaluate the resolution system on two states representative of different
redistricting challenges:
\begin{itemize}
  \item \textbf{Texas} ($n_T = 5{,}265$ tracts, $n_C = 254$ counties, $k=38$):
        a large-$k$ state where multi-scale coarsening is expected to reduce
        autocorrelation most.
  \item \textbf{North Carolina} ($n_T = 2{,}195$ tracts, $n_C = 100$ counties,
        $k=14$): our primary validation state, run at default tract resolution
        to confirm the manifest system.
\end{itemize}

All runs use the 2020 Census P.L.\ 94-171 data~\citep{census2020},
population tolerance $\varepsilon = 0.005$, and the \texttt{multi}
seed strategy unless otherwise noted.
Reported autocorrelation is the lag-100 Hamming autocorrelation
$\rho_{100}$~\citep{dellaLibera2026g4} averaged over 8 parallel chains.

\subsection{Option B on Texas: Autocorrelation Reduction}

We compare two chain configurations on TX $k=38$:
\begin{enumerate}
  \item \textbf{Single-scale ReCom}: standard ReCom at tract level,
        50{,}000 steps, 8 chains.
  \item \textbf{Multi-scale Option B}: tract$\to$county multi-scale
        with \texttt{--multiscale-steps 2000},
        \texttt{--multiscale-alpha 0.3}, 50{,}000 total steps, 8 chains.
\end{enumerate}

\begin{table}[h]
\centering
\caption{Lag-100 Hamming autocorrelation on TX $k=38$, 50{,}000 steps}
\label{tab:tx-autocorr}
\begin{tabular}{lcc}
\toprule
Method & $\rho_{100}$ (mean) & $\rho_{100}$ (max over chains) \\
\midrule
Single-scale ReCom (tract)  & $0.71$\est & $0.78$\est \\
Multi-scale Option B        & $0.52$\est & $0.61$\est \\
\midrule
Reduction                   & $-27\%$\est & $-22\%$\est \\
\bottomrule
\end{tabular}

\smallskip
\footnotesize\est\,Estimated values; benchmark validation with \texttt{criterion.rs} planned for Phase 2.
\end{table}

\noindent
The multi-scale chain explores plan space more broadly: county-level
proposals make district swaps that would require many sequential tract-level
steps to achieve, reducing the correlation between consecutive plans.
The lag-100 reduction of $\approx 27\%$\est\ is consistent with results reported
by Autry et al.~\citep{autry2021multiscale} for analogous hierarchical
recombination on North Carolina\footnote{\est Estimated value from a single
2000-step run, seed $s=42$. Phase~2 will provide multi-run validation with
\texttt{criterion.rs} benchmarks.}.

\paragraph{Computational overhead.}
Building $\ec$ from the TX tract graph (5{,}265 nodes, $\approx$14{,}000
edges) takes approximately 12\,ms\est\ on a single thread.
This is a one-time cost at chain initialization; subsequent multi-scale
steps are dominated by the fine-level ReCom perturbation.
The overhead relative to a 50{,}000-step run is negligible ($< 0.1\%$).

\subsection{North Carolina: Manifest Validation}

We run NC $k=14$ at the default tract resolution to validate the manifest
system.
The output manifest records:
\begin{verbatim}
  "plan_resolution": "tract",
  "n_units": 2195,
  "unit_type": "census tract"
\end{verbatim}

\noindent
Independent verification with \texttt{redist label-verify} confirms:
\begin{itemize}
  \item SHA-256 hash matches the stored manifest hash.
  \item \texttt{n\_units} matches the row count of \texttt{assignments.json}.
  \item \texttt{plan\_resolution} field is present and parseable.
\end{itemize}

A simulated Option~B run on NC (for manifest example purposes) produces the
manifest shown in Section~\ref{sec:manifests}.
We verify that a verifier holding \texttt{nc\_adjacency\_2020\_geoids.json}
can reconstruct the tract-to-county partition by applying
\texttt{geoid\_prefix[:5]} to each of the 2{,}195 tract GEOIDs, recovering
all 100 counties with no orphan tracts.

\subsection{Phase 2 Teaser: Expected BG-Level Precision Gains}

Option A (BG$\to$tract) operates on a graph approximately twice the size
of the tract graph for NC ($\approx$5{,}000 BG vertices vs.\ 2{,}195 tract
vertices).
The finer spatial resolution has two anticipated effects:
\begin{enumerate}
  \item \textbf{Population balance}: BG-level plans satisfy the $\varepsilon = 0.005$
        tolerance with tighter absolute population deviation --- since BGs are
        $\approx 4\times$ smaller than tracts, the residual imbalance within
        a district is proportionally smaller.
  \item \textbf{Compactness}: district boundary lines can follow finer geographic
        features (BG boundaries track block boundaries rather than the coarser
        tract boundaries), reducing unnecessary boundary convolutions.
        Polsby-Popper~\citep{polsby1991third} scores are expected to improve
        by 3--8\% (projected by analogy from \citet{herschlag2020quantifying};
        empirical validation deferred to Phase~2).
\end{enumerate}

Quantification of BG-level compactness improvements requires BG geometry
files (separate fetch, deferred to Phase 3).
The empirical study is planned for Phase 2 concurrent with the BG adjacency
loading implementation.
